% !TeX program = lualatex
% !TeX encoding = utf8
% !TeX spellcheck = uk_UA
% !BIB program = bibtex8

\documentclass[12pt]{scrbook}
\usepackage{MyBook}




\begin{document}


\problem{}{1}{
Виразіть температуру поверхні Сонця в енергетичних одиницях: еВ, ерг, Дж.
}

\problem{}{2}{
Довгим соленоїдом, що має $n$ витків на одиницю довжини, тече струм $I$.
Знайти тиск, що діє на бічну поверхню соленоїда. \emph{Вказівка: викладки проводьте в гаусовій системі одиниць.}
}

\problem{}{3}{
Чи може існувати хвиля вигдяду:
\begin{equation*}
    \vec{E} = \vec{e}_x E_0 e^{i(\omega t - kx)} ?
\end{equation*}
}

\problem{}{4}{
    Плоска ехвиля з вектором напруженості електричного поля $\vec{E}$, що лежить в площині падіння потрапляє із середовища з показником заломлення $n_1$ по нормалі на плоску границю розділу з середовищем, що має показник заломлення $n_2$. Отримати  вирази для коефіцієнтів відбивання та пропускання хвиль.  \emph{Вказівка: для виведення коефіцієнтів користуйтесь експоненційною формою запису плоскої електромагнітної хвилі.}
}
\end{document}
